\documentclass[12pt]{article}

\usepackage[spanish]{babel}
\usepackage[utf8]{inputenc}
\usepackage{amsmath, amssymb}
\usepackage{geometry}
\usepackage{graphicx}
\usepackage{float}
\usepackage{booktabs}

\geometry{letterpaper, margin=2.5cm}

\title{\textbf{Actividad 3.2 \\ Análisis de Procesos Estocásticos \\ Ejemplo 8: Servidores Cloud}}
\author{Emilio Izquierdo Montero \\ 8SA}
\date{23 de Abril 2026}

\begin{document}

\maketitle

%------------------------------------------------
\section{Introducción}

En este trabajo se analiza un proceso estocástico multivariante utilizando datos de dos variables observadas: uso de CPU y uso de memoria. El objetivo es identificar estados ocultos del sistema mediante Modelos de Mezcla Gaussiana (GMM) y analizar su comportamiento.

%------------------------------------------------
\section{Descripción del Problema}

Las variables analizadas son:

\begin{itemize}
\item $X(t)$: Uso de CPU (\%)
\item $Y(t)$: Uso de memoria RAM (\%)
\end{itemize}

Se busca identificar estados ocultos como:

\begin{itemize}
\item Bajo uso (Idle)
\item Uso intensivo de CPU
\item Uso intensivo de memoria
\item Sobrecarga
\end{itemize}

%------------------------------------------------
\section{Exploración de Datos}

Se realizó un análisis visual inicial mediante un gráfico de dispersión.

\begin{figure}[H]
\centering
\includegraphics[width=0.7\textwidth]{/home/emilio/proyectos/adan1/figura1.png}
\caption{Distribución de las variables (Variable\_X vs Variable\_Y)}
\end{figure}

\textbf{Análisis:} Se observan agrupaciones en el espacio de fases, lo que sugiere la existencia de múltiples estados en el sistema.

%------------------------------------------------
\section{Selección del Número de Estados}

Se utilizó el criterio BIC para determinar el número óptimo de clusters.

\begin{figure}[H]
\centering
\includegraphics[width=0.7\textwidth]{/home/emilio/proyectos/adan1/clusters.png}
\caption{Selección del número de clusters mediante BIC}
\end{figure}

\textbf{Resultado:} El número óptimo de estados es $K = X$ (sustituir por resultado real).

%------------------------------------------------
\section{Modelado con GMM}

El modelo utilizado se define como:

\begin{equation}
p(x) = \sum_{k=1}^{K} \pi_k \mathcal{N}(x | \mu_k, \Sigma_k)
\end{equation}

Se entrenó un modelo GMM con $K$ componentes.

\begin{figure}[H]
\centering
\includegraphics[width=0.7\textwidth]{/home/emilio/proyectos/adan1/seleccionk.png}
\caption{Clusters identificados por el modelo GMM}
\end{figure}

%------------------------------------------------
\section{Resultados}

\subsection{Agrupación en el espacio de fases}

Los datos se agrupan en regiones diferenciadas, indicando distintos estados operativos.

\subsection{Parámetros del modelo}

Cada estado se caracteriza por su media y covarianza.

\begin{itemize}
\item Las medias $\mu$ representan el centro de cada cluster
\item Las covarianzas $\Sigma$ describen la dispersión
\item Los pesos $\pi$ indican la proporción de datos en cada estado
\end{itemize}

\textbf{(Opcional)} Puedes incluir aquí una tabla con resultados:

\begin{table}[H]
\centering
\begin{tabular}{ccc}
\toprule
Estado & $\mu_X$ & $\mu_Y$ \\
\midrule
0 & ... & ... \\
1 & ... & ... \\
2 & ... & ... \\
3 & ... & ... \\
\bottomrule
\end{tabular}
\caption{Medias de los clusters}
\end{table}

%------------------------------------------------
\section{Modelo HMM (Opcional)}

Para considerar la dinámica temporal, se utilizó un modelo oculto de Markov.

\begin{figure}[H]
\centering
\includegraphics[width=0.7\textwidth]{/home/emilio/proyectos/adan1/estadis.png}
\caption{Estados ocultos en el tiempo}
\end{figure}

\subsection{Matriz de transición}

\[
A =
\begin{bmatrix}
P_{11} & P_{12} & \cdots \\
P_{21} & P_{22} & \cdots \\
\vdots & \vdots & \ddots
\end{bmatrix}
\]

\textbf{Interpretación:} Esta matriz describe la probabilidad de transición entre estados.

%------------------------------------------------
\section{Análisis e Interpretación}

\subsection{Agrupación}
Las observaciones se agrupan en regiones que corresponden a diferentes comportamientos del sistema.

\subsection{Características de los estados}
Cada estado presenta valores característicos de CPU y memoria, lo que permite interpretarlos como distintos modos de operación.

\subsection{Alineación con estados reales}
Los clusters identificados presentan similitudes con los estados teóricos, aunque existe cierto solapamiento entre ellos.

\subsection{Aplicaciones}
Estas técnicas permiten:

\begin{itemize}
\item Detectar sobrecarga del sistema
\item Optimizar recursos
\item Implementar autoescalado
\item Identificar anomalías
\end{itemize}

\subsection{Probabilidades de transición}
El modelo HMM permite calcular la probabilidad de cambio entre estados, útil para predicción y toma de decisiones.

%------------------------------------------------
\section{Conclusión}

Los modelos de mezcla gaussiana permiten identificar estados ocultos en sistemas complejos a partir de datos observados. La incorporación de modelos HMM permite además analizar la evolución temporal, mejorando la comprensión del sistema.

\end{document}
